% main.tex
% Minimal, monochrome penetration-testing report template
% Compile: pdflatex main.tex (run twice for the table of contents)
% Optional: latexmk -pdf main.tex
%
% Dynamic content can later be populated by Python/Jinja2.
% Keep report data separate from this presentation template.

\documentclass[11pt,a4paper]{report}

% ---------------------------------------------------------------------------
% Packages
% ---------------------------------------------------------------------------
\usepackage[utf8]{inputenc}
\usepackage[T1]{fontenc}
\usepackage{lmodern}
\usepackage[a4paper,margin=25mm,headheight=15pt]{geometry}
\usepackage{microtype}
\usepackage{graphicx}
\usepackage{float}
\usepackage{booktabs}
\usepackage{longtable}
\usepackage{tabularx}
\usepackage{array}
\usepackage{enumitem}
\usepackage{fancyhdr}
\usepackage{titlesec}
\usepackage{caption}
\usepackage{hyperref}
\usepackage{xcolor}
\usepackage{listings}
\usepackage{verbatim}
\usepackage{datetime2}

% ---------------------------------------------------------------------------
% Neutral style variables
% These are intentionally centralized so they can be overridden later.
% No blue/teal colors are used.
% ---------------------------------------------------------------------------
\definecolor{ReportBlack}{HTML}{000000}
\definecolor{ReportDarkGray}{HTML}{333333}
\definecolor{ReportMidGray}{HTML}{666666}
\definecolor{ReportLightGray}{HTML}{EEEEEE}
\definecolor{ReportRuleGray}{HTML}{BBBBBB}

\hypersetup{
    colorlinks=true,
    linkcolor=ReportBlack,
    urlcolor=ReportBlack,
    citecolor=ReportBlack,
    pdfborder={0 0 0}
}

% ---------------------------------------------------------------------------
% Report metadata / replaceable values
% ---------------------------------------------------------------------------
\newcommand{\ReportTitle}{Penetration Testing Report}
\newcommand{\ClientName}{[CLIENT / ORGANIZATION]}
\newcommand{\ProjectName}{[PROJECT NAME]}
\newcommand{\PreparedBy}{[AUTHOR / SECURITY TEAM]}
\newcommand{\ReportDate}{[REPORT DATE]}
\newcommand{\AssessmentStart}{[START DATE]}
\newcommand{\AssessmentEnd}{[END DATE]}
\newcommand{\TargetEnvironment}{[TARGET ENVIRONMENT]}
\newcommand{\TargetApplication}{VAmPI -- Vulnerable API}
\newcommand{\ReportVersion}{1.0}
\newcommand{\Classification}{CONFIDENTIAL}
\newcommand{\OverallRisk}{[OVERALL RISK]}
\newcommand{\TestingType}{[BLACK BOX / GREY BOX / WHITE BOX]}
\newcommand{\CommitSHA}{[GIT COMMIT SHA]}
\newcommand{\PipelineID}{[GITLAB PIPELINE ID]}
\newcommand{\PipelineURL}{[GITLAB PIPELINE URL]}
\newcommand{\RepositoryURL}{[REPOSITORY URL]}
\newcommand{\BranchName}{[BRANCH]}
\newcommand{\ReportStatus}{[DRAFT / FINAL]}

% Severity totals: populate from normalized scanner/test data.
\newcommand{\CriticalCount}{[N]}
\newcommand{\HighCount}{[N]}
\newcommand{\MediumCount}{[N]}
\newcommand{\LowCount}{[N]}
\newcommand{\InfoCount}{[N]}
\newcommand{\TotalFindings}{[N]}

% ---------------------------------------------------------------------------
% Formatting
% ---------------------------------------------------------------------------
\setlength{\parindent}{0pt}
\setlength{\parskip}{6pt}
\renewcommand{\arraystretch}{1.2}
\setlist[itemize]{leftmargin=*,itemsep=2pt,topsep=3pt}
\setlist[enumerate]{leftmargin=*,itemsep=2pt,topsep=3pt}

\titleformat{\chapter}
  {\normalfont\huge\bfseries\color{ReportBlack}}
  {\thechapter}{1em}{}
\titlespacing*{\chapter}{0pt}{-10pt}{18pt}

\titleformat{\section}
  {\normalfont\Large\bfseries\color{ReportBlack}}
  {\thesection}{1em}{}
\titlespacing*{\section}{0pt}{16pt}{7pt}

\titleformat{\subsection}
  {\normalfont\large\bfseries\color{ReportDarkGray}}
  {\thesubsection}{1em}{}
\titlespacing*{\subsection}{0pt}{12pt}{5pt}

\captionsetup{
    font=small,
    labelfont=bf,
    textfont=normalfont
}

\pagestyle{fancy}
\fancyhf{}
\fancyhead[L]{\small\ReportTitle}
\fancyhead[R]{\small\ProjectName}
\fancyfoot[L]{\small\Classification}
\fancyfoot[R]{\small Page \thepage}
\renewcommand{\headrulewidth}{0.4pt}
\renewcommand{\footrulewidth}{0.4pt}

\newcolumntype{Y}{>{\raggedright\arraybackslash}X}
\newcommand{\placeholder}[1]{\textbf{[#1]}}
\newcommand{\findingfield}[2]{\textbf{#1}\\#2\par}
\newcommand{\severity}[1]{\textbf{#1}}

% Listing style for small evidence/configuration excerpts.
\lstset{
    basicstyle=\ttfamily\small,
    breaklines=true,
    columns=fullflexible,
    frame=single,
    rulecolor=\color{ReportRuleGray},
    backgroundcolor=\color{white},
    numbers=left,
    numberstyle=\tiny\color{ReportMidGray},
    showstringspaces=false,
    tabsize=2
}

% ---------------------------------------------------------------------------
% Document
% ---------------------------------------------------------------------------
\begin{document}

% ---------------------------------------------------------------------------
% Cover page
% ---------------------------------------------------------------------------
\begin{titlepage}
    \thispagestyle{empty}
    \vspace*{18mm}

    {\Large\bfseries \ReportTitle\par}
    \vspace{4mm}
    {\large \ProjectName\par}

    \vspace{18mm}
    \hrule height 0.8pt
    \vspace{10mm}

    \begin{tabularx}{\textwidth}{@{}>{\bfseries}p{43mm}Y@{}}
        Prepared for: & \ClientName \\
        Prepared by: & \PreparedBy \\
        Application: & \TargetApplication \\
        Environment: & \TargetEnvironment \\
        Assessment type: & \TestingType \\
        Assessment dates: & \AssessmentStart\ -- \AssessmentEnd \\
        Report date: & \ReportDate \\
        Report version: & \ReportVersion \\
        Status: & \ReportStatus \\
    \end{tabularx}

    \vfill
    \hrule height 0.4pt
    \vspace{4mm}
    {\small\bfseries \Classification\par}
    {\small This document contains security assessment information intended
    only for the authorized recipient. Do not distribute without permission.\par}
\end{titlepage}

\pagenumbering{roman}

% ---------------------------------------------------------------------------
% Document control
% ---------------------------------------------------------------------------
\chapter*{Document Control}
\addcontentsline{toc}{chapter}{Document Control}

\begin{tabularx}{\textwidth}{@{}p{43mm}Y@{}}
    \toprule
    \textbf{Field} & \textbf{Value} \\
    \midrule
    Document title & \ReportTitle \\
    Project & \ProjectName \\
    Client & \ClientName \\
    Version & \ReportVersion \\
    Report status & \ReportStatus \\
    Classification & \Classification \\
    Author & \PreparedBy \\
    Report date & \ReportDate \\
    \bottomrule
\end{tabularx}

\vspace{8mm}
\textbf{Revision history}

\begin{tabularx}{\textwidth}{@{}p{20mm}p{32mm}p{35mm}Y@{}}
    \toprule
    \textbf{Version} & \textbf{Date} & \textbf{Author} & \textbf{Description} \\
    \midrule
    1.0 & \ReportDate & \PreparedBy & Initial report \\
    \placeholder{VERSION} & \placeholder{DATE} & \placeholder{AUTHOR} &
    \placeholder{CHANGE DESCRIPTION} \\
    \bottomrule
\end{tabularx}

\tableofcontents
\listoftables
\listoffigures
\clearpage

\pagenumbering{arabic}

% ---------------------------------------------------------------------------
% 1. Executive Summary
% ---------------------------------------------------------------------------
\chapter{Executive Summary}

\section{Assessment Overview}
This report documents the results of a security assessment of
\TargetApplication, conducted for \ClientName. The assessment focused on the
application and its supporting CI/CD security controls in the
\TargetEnvironment\ environment.

The assessment was performed from \AssessmentStart\ to \AssessmentEnd.
The testing approach was \TestingType.

\section{Objectives}
The objectives were to:
\begin{itemize}
    \item Assess the security posture of the deliberately vulnerable VAmPI API.
    \item Evaluate security checks integrated into the GitLab CI/CD pipeline.
    \item Identify exposed secrets, application security weaknesses, vulnerable
          dependencies and container image risks.
    \item Review software bill of materials (SBOM) generation and vulnerability
          analysis.
    \item Review image signing and signature verification controls.
    \item Provide actionable remediation recommendations.
\end{itemize}

\section{Overall Assessment}
\textbf{Overall risk rating:} \OverallRisk

\textbf{Summary:} \placeholder{WRITE A SHORT, EVIDENCE-BASED SUMMARY OF THE
MOST IMPORTANT RESULTS, BUSINESS/TECHNICAL IMPACT, AND OVERALL PIPELINE POSTURE.}

\section{Findings Summary}
\begin{table}[H]
\centering
\caption{Findings by severity}
\begin{tabularx}{0.85\textwidth}{@{}Y r@{}}
    \toprule
    \textbf{Severity} & \textbf{Count} \\
    \midrule
    Critical & \CriticalCount \\
    High & \HighCount \\
    Medium & \MediumCount \\
    Low & \LowCount \\
    Informational & \InfoCount \\
    \midrule
    \textbf{Total} & \textbf{\TotalFindings} \\
    \bottomrule
\end{tabularx}
\end{table}

% ---------------------------------------------------------------------------
% 2. Scope and Methodology
% ---------------------------------------------------------------------------
\chapter{Scope and Methodology}

\section{Scope}
\begin{tabularx}{\textwidth}{@{}p{43mm}Y@{}}
    \toprule
    \textbf{Item} & \textbf{Description} \\
    \midrule
    Target application & \TargetApplication \\
    Target environment & \TargetEnvironment \\
    Repository & \RepositoryURL \\
    Branch & \BranchName \\
    Commit & \CommitSHA \\
    GitLab pipeline ID & \PipelineID \\
    Pipeline URL & \PipelineURL \\
    Assessment type & \TestingType \\
    \bottomrule
\end{tabularx}

\section{In-Scope Components}
\begin{itemize}
    \item VAmPI application and exposed API endpoints.
    \item Application source code included in the assessed repository.
    \item GitLab CI/CD configuration and security pipeline stages.
    \item Container image built by the assessed pipeline.
    \item Dependency and package inventory represented by the generated SBOM.
    \item Container image vulnerability scanning and signing/verification steps.
\end{itemize}

\section{Out-of-Scope Components}
\begin{itemize}
    \item \placeholder{LIST EXCLUDED HOSTS, SERVICES, NETWORKS OR COMPONENTS}
    \item \placeholder{IDENTIFY ANY THIRD-PARTY SYSTEMS NOT AUTHORIZED FOR TESTING}
\end{itemize}

\section{Rules of Engagement and Limitations}
\placeholder{DOCUMENT AUTHORIZATION, TESTING WINDOWS, RATE LIMITS,
PROHIBITED ACTIONS, TEST ACCOUNTS, ENVIRONMENTAL LIMITATIONS AND ANY
INTERRUPTIONS OR CONSTRAINTS.}

\section{Methodology}
Describe the methodology actually followed. Relevant references may include:
\begin{itemize}
    \item OWASP Web Security Testing Guide.
    \item OWASP API Security Top 10.
    \item Penetration Testing Execution Standard (PTES).
    \item NIST SP 800-115, Technical Guide to Information Security Testing
          and Assessment.
\end{itemize}
\textit{Only list a methodology or reference as applied when it was actually
used during the assessment.}

\section{Tools and Techniques}
\begin{longtable}{@{}p{35mm}p{38mm}p{75mm}@{}}
    \caption{Assessment tools}\\
    \toprule
    \textbf{Tool} & \textbf{Purpose} & \textbf{Version / Configuration} \\
    \midrule
    \endfirsthead
    \toprule
    \textbf{Tool} & \textbf{Purpose} & \textbf{Version / Configuration} \\
    \midrule
    \endhead
    Gitleaks & Secret detection & \placeholder{VERSION / CONFIGURATION} \\
    Semgrep & Static application security testing & \placeholder{VERSION / RULESETS} \\
    Syft & SBOM generation & \placeholder{VERSION / OUTPUT FORMAT} \\
    Trivy & Container / dependency scanning & \placeholder{VERSION / CONFIGURATION} \\
    Grype & SBOM vulnerability analysis & \placeholder{VERSION / CONFIGURATION} \\
    Cosign & Image signing and verification & \placeholder{VERSION / VERIFICATION MODE} \\
    \placeholder{OTHER TOOL} & \placeholder{PURPOSE} & \placeholder{VERSION} \\
    \bottomrule
\end{longtable}

% ---------------------------------------------------------------------------
% 3. CI/CD Security Architecture
% ---------------------------------------------------------------------------
\chapter{CI/CD Security Architecture}

\section{Pipeline Overview}
\placeholder{DESCRIBE THE PIPELINE FLOW, TRIGGERS, RUNNERS, REGISTRIES,
ARTIFACT STORAGE AND DEPLOYMENT TARGET. INSERT AN ARCHITECTURE DIAGRAM
IF AVAILABLE.}

% Example:
% \begin{figure}[H]
%   \centering
%   \includegraphics[width=0.9\textwidth]{evidence/pipeline-architecture.png}
%   \caption{CI/CD security pipeline architecture}
% \end{figure}

\section{Security Controls by Pipeline Stage}
\begin{longtable}{@{}p{29mm}p{37mm}p{79mm}@{}}
    \caption{Pipeline security control mapping}\\
    \toprule
    \textbf{Stage} & \textbf{Control} & \textbf{Expected Evidence / Result} \\
    \midrule
    \endfirsthead
    \toprule
    \textbf{Stage} & \textbf{Control} & \textbf{Expected Evidence / Result} \\
    \midrule
    \endhead
    Source validation & Gitleaks & Secret scan report and detected-secret review \\
    SAST & Semgrep & Static analysis findings and rule identifiers \\
    SBOM generation & Syft & CycloneDX or SPDX inventory artifact \\
    Image scanning & Trivy & Image/package vulnerability report \\
    SBOM analysis & Grype & Vulnerability matches against SBOM components \\
    Image signing & Cosign & Signature and signing identity evidence \\
    Verification & Cosign & Verification result and policy outcome \\
    \bottomrule
\end{longtable}

\section{Pipeline Execution Details}
\begin{tabularx}{\textwidth}{@{}p{43mm}Y@{}}
    \toprule
    \textbf{Field} & \textbf{Value} \\
    \midrule
    GitLab project & \placeholder{PROJECT PATH} \\
    Pipeline ID & \PipelineID \\
    Pipeline URL & \PipelineURL \\
    Branch / tag & \BranchName \\
    Commit SHA & \CommitSHA \\
    Execution date & \placeholder{DATE AND TIME WITH TIME ZONE} \\
    Pipeline result & \placeholder{PASSED / FAILED / PARTIAL} \\
    Runner / executor & \placeholder{RUNNER DETAILS} \\
    Image reference & \placeholder{IMAGE NAME AND TAG / DIGEST} \\
    \bottomrule
\end{tabularx}

% ---------------------------------------------------------------------------
% 4. Findings Summary
% ---------------------------------------------------------------------------
\chapter{Findings Summary}

The following table summarizes the security findings identified and validated
during the assessment. Findings should be assigned a status that reflects their
actual validation state.

\begin{longtable}{@{}p{16mm}p{65mm}p{25mm}p{27mm}p{25mm}@{}}
    \caption{Summary of security findings}\\
    \toprule
    \textbf{ID} & \textbf{Finding} & \textbf{Severity} &
    \textbf{Status} & \textbf{Source} \\
    \midrule
    \endfirsthead
    \toprule
    \textbf{ID} & \textbf{Finding} & \textbf{Severity} &
    \textbf{Status} & \textbf{Source} \\
    \midrule
    \endhead
    F-001 & \placeholder{FINDING TITLE} & \placeholder{SEVERITY} &
    \placeholder{STATUS} & \placeholder{TOOL / TEST} \\
    F-002 & \placeholder{FINDING TITLE} & \placeholder{SEVERITY} &
    \placeholder{STATUS} & \placeholder{TOOL / TEST} \\
    \placeholder{F-ID} & \placeholder{FINDING TITLE} &
    \placeholder{SEVERITY} & \placeholder{STATUS} &
    \placeholder{TOOL / TEST} \\
    \bottomrule
\end{longtable}

\textbf{Status definitions:}
\begin{itemize}
    \item \textbf{Detected:} Reported by a tool or test and awaiting validation.
    \item \textbf{Validated:} Reviewed and confirmed as a security issue.
    \item \textbf{False Positive:} Reviewed and determined not to represent a
          valid issue in the assessed context.
    \item \textbf{Accepted Risk:} Formally accepted by the authorized risk owner.
    \item \textbf{Remediated:} Fix implemented and verified through retesting.
\end{itemize}

% ---------------------------------------------------------------------------
% 5. Detailed Findings
% ---------------------------------------------------------------------------
\chapter{Detailed Findings}

% Duplicate the following section for each finding.
\section{F-001 -- \placeholder{FINDING TITLE}}

\begin{tabularx}{\textwidth}{@{}p{43mm}Y@{}}
    \toprule
    \textbf{Attribute} & \textbf{Value} \\
    \midrule
    Finding ID & F-001 \\
    Title & \placeholder{FINDING TITLE} \\
    Category & \placeholder{SECRETS / SAST / DEPENDENCY / CONTAINER / API / SUPPLY CHAIN} \\
    Severity & \placeholder{CRITICAL / HIGH / MEDIUM / LOW / INFORMATIONAL} \\
    CVSS version & \placeholder{CVSS 3.1 OR 4.0, IF APPLICABLE} \\
    CVSS score & \placeholder{SCORE} \\
    CVSS vector & \placeholder{VECTOR STRING} \\
    Affected asset & \placeholder{FILE / ENDPOINT / PACKAGE / IMAGE / PIPELINE JOB} \\
    Source tool / test & \placeholder{TOOL OR MANUAL TEST} \\
    Status & \placeholder{DETECTED / VALIDATED / FALSE POSITIVE / ACCEPTED RISK / REMEDIATED} \\
    References & \placeholder{CVE / CWE / OWASP / ADVISORY / RULE ID} \\
    \bottomrule
\end{tabularx}

\subsection{Description}
\placeholder{DESCRIBE THE ISSUE, ITS ROOT CAUSE, AND THE CONDITIONS UNDER
WHICH IT WAS IDENTIFIED. DISTINGUISH TOOL OUTPUT FROM MANUAL VALIDATION.}

\subsection{Evidence}
\placeholder{INSERT REDACTED EVIDENCE, REPRODUCTION STEPS, SCREENSHOTS,
RELEVANT LOG EXCERPTS OR SCANNER OUTPUT. DO NOT INCLUDE LIVE SECRETS,
PASSWORDS, TOKENS OR OTHER SENSITIVE VALUES.}

% Example evidence image:
% \begin{figure}[H]
%   \centering
%   \includegraphics[width=0.9\textwidth]{evidence/f-001.png}
%   \caption{Evidence for finding F-001}
% \end{figure}

% Example evidence excerpt:
% \begin{lstlisting}
% [Insert sanitized tool output or command result]
% \end{lstlisting}

\subsection{Impact}
\placeholder{EXPLAIN THE TECHNICAL IMPACT AND, WHERE RELEVANT, THE
POTENTIAL OPERATIONAL OR BUSINESS CONSEQUENCES.}

\subsection{Likelihood and Exploitability}
\placeholder{DESCRIBE REQUIRED ACCESS, ATTACK COMPLEXITY, USER INTERACTION,
EXPOSURE AND OTHER RELEVANT CONDITIONS.}

\subsection{Remediation}
\begin{enumerate}
    \item \placeholder{IMMEDIATE CONTAINMENT OR MITIGATION}
    \item \placeholder{ROOT-CAUSE FIX}
    \item \placeholder{PREVENTIVE CONTROL OR PIPELINE IMPROVEMENT}
    \item \placeholder{RETEST / VERIFICATION PROCEDURE}
\end{enumerate}

\subsection{Retest Result}
\placeholder{NOT RETESTED / FIX VERIFIED / ISSUE STILL PRESENT / PARTIALLY FIXED.
INCLUDE RETEST DATE, COMMIT OR PIPELINE REFERENCE WHEN AVAILABLE.}

% ---------------------------------------------------------------------------
% Copy the complete F-001 section above for F-002, F-003, etc.
% A future Python/Jinja2 generator can render one section per normalized
% finding object and populate the summary table and severity totals.
% ---------------------------------------------------------------------------

% ---------------------------------------------------------------------------
% 6. Risk Rating Method
% ---------------------------------------------------------------------------
\chapter{Risk Rating Method}

\section{Severity Classification}
If CVSS is used, state the CVSS version and scoring assumptions. The following
CVSS 3.1 bands are provided as a reference.

\begin{table}[H]
\centering
\caption{CVSS 3.1 severity bands}
\begin{tabularx}{0.85\textwidth}{@{}Y c@{}}
    \toprule
    \textbf{Severity} & \textbf{CVSS score} \\
    \midrule
    Critical & 9.0 -- 10.0 \\
    High & 7.0 -- 8.9 \\
    Medium & 4.0 -- 6.9 \\
    Low & 0.1 -- 3.9 \\
    Informational & 0.0 / Not scored \\
    \bottomrule
\end{tabularx}
\end{table}

CVSS is a technical severity measure and should not be treated as a complete
business-risk assessment. Consider exposure, exploitability, asset criticality,
existing controls and the assessment context when assigning the overall risk.

\section{Risk Acceptance}
\placeholder{DOCUMENT THE RISK ACCEPTANCE PROCESS, APPROVAL AUTHORITY,
EXPIRATION DATE AND REQUIRED COMPENSATING CONTROLS, IF APPLICABLE.}

% ---------------------------------------------------------------------------
% 7. Remediation Recommendations
% ---------------------------------------------------------------------------
\chapter{Remediation Recommendations}

\section{Prioritized Remediation Plan}
\begin{longtable}{@{}p{22mm}p{65mm}p{28mm}p{35mm}@{}}
    \caption{Remediation plan}\\
    \toprule
    \textbf{Priority} & \textbf{Action} & \textbf{Owner} & \textbf{Target date} \\
    \midrule
    \endfirsthead
    \toprule
    \textbf{Priority} & \textbf{Action} & \textbf{Owner} & \textbf{Target date} \\
    \midrule
    \endhead
    Immediate & \placeholder{ACTION} & \placeholder{OWNER} & \placeholder{DATE} \\
    Short term & \placeholder{ACTION} & \placeholder{OWNER} & \placeholder{DATE} \\
    Medium term & \placeholder{ACTION} & \placeholder{OWNER} & \placeholder{DATE} \\
    Long term & \placeholder{ACTION} & \placeholder{OWNER} & \placeholder{DATE} \\
    \bottomrule
\end{longtable}

\section{Recommended CI/CD Improvements}
Potential improvement areas, where supported by the assessment results:
\begin{itemize}
    \item Define failure thresholds for high-confidence, validated security findings.
    \item Protect and restrict access to CI/CD variables and signing credentials.
    \item Preserve scan reports and SBOM artifacts with appropriate access controls.
    \item Pin critical build dependencies and container images to immutable digests.
    \item Verify image signatures before deployment, using an explicit trust policy.
    \item Define exception handling, expiry dates and approval for accepted risks.
    \item Re-run relevant security checks after remediation.
\end{itemize}

\section{Retesting}
\placeholder{DESCRIBE THE RETEST SCOPE, RESPONSIBLE PARTY, EXPECTED EVIDENCE
AND HOW REMEDIATION WILL BE CONFIRMED.}

% ---------------------------------------------------------------------------
% 8. Conclusion
% ---------------------------------------------------------------------------
\chapter{Conclusion}

\placeholder{SUMMARIZE THE ASSESSMENT OUTCOME, MOST IMPORTANT RISKS,
THE EFFECTIVENESS OF THE TESTED CONTROLS AND THE PRIORITY NEXT STEPS.
ENSURE THAT THE CONCLUSION IS CONSISTENT WITH THE EVIDENCE AND FINDINGS
DOCUMENTED IN THIS REPORT.}

% ---------------------------------------------------------------------------
% Appendices
% ---------------------------------------------------------------------------
\appendix

\chapter{Pipeline Execution Evidence}
\placeholder{INCLUDE RELEVANT PIPELINE JOB LOG EXCERPTS, JOB STATUS,
ARTIFACT REFERENCES AND EXECUTION METADATA. REDACT SECRETS AND SENSITIVE DATA.}

\chapter{Scanner Reports and SBOM}
\placeholder{REFERENCE THE ORIGINAL JSON / SARIF / CYCLONEDX / SPDX
ARTIFACTS. INCLUDE SELECTED, SANITIZED EXTRACTS ONLY WHEN THEY HELP
EXPLAIN A FINDING.}

\chapter{Screenshots and Supporting Evidence}
% Add evidence images as needed:
% \begin{figure}[H]
%   \centering
%   \includegraphics[width=\textwidth]{evidence/example.png}
%   \caption{Description of evidence}
% \end{figure}
\placeholder{ADD SCREENSHOTS, DIAGRAMS OR OTHER SUPPORTING MATERIAL.
USE STABLE FILENAMES AND REFERENCE THEM FROM THE RELEVANT FINDINGS.}

\chapter{Glossary}
\begin{longtable}{@{}p{35mm}p{115mm}@{}}
    \toprule
    \textbf{Term} & \textbf{Definition} \\
    \midrule
    API & Application Programming Interface \\
    CI/CD & Continuous Integration / Continuous Delivery or Deployment \\
    CVSS & Common Vulnerability Scoring System \\
    SAST & Static Application Security Testing \\
    SBOM & Software Bill of Materials \\
    \placeholder{TERM} & \placeholder{DEFINITION} \\
    \bottomrule
\end{longtable}

\chapter{Assumptions and Limitations}
\placeholder{RECORD RELEVANT ASSUMPTIONS, TESTING LIMITATIONS, ENVIRONMENTAL
CONSTRAINTS, UNTESTED COMPONENTS AND ANY OTHER QUALIFICATIONS THAT AFFECT
INTERPRETATION OF THE RESULTS.}

\end{document}
